\section*{1. Means, covariance, and a fixed design}
For a random column $Y\in\R^n$ with finite second moments, define
\[
\mu=\E Y,\qquad \Cov(Y)=\E[(Y-\mu)(Y-\mu)^\top].
\]
Expectation of a vector or matrix is taken entry by entry. A covariance
matrix is symmetric and positive semidefinite, since
$a^\top\Cov(Y)a=\E[(a^\top(Y-\mu))^2]\geqslant0$.
For deterministic matrices $A,B$ and a deterministic column $c$,
\[
\E(AY+c)=A\mu+c,\quad \Cov(AY)=A\Cov(Y)A^\top,\quad
\Cov(AY,BY)=A\Cov(Y)B^\top.
\]
For example, the last identity follows by substituting
$AY-\E(AY)=A(Y-\mu)$ and $BY-\E(BY)=B(Y-\mu)$ in the definition of
cross-covariance and moving the fixed matrices outside the expectation.

\begin{proposition}[Expected quadratic form]\label{prop:11:quadratic}
For a fixed $n\times n$ matrix $A$,
\[
\E(Y^\top AY)=\tr(A\Cov(Y))+\mu^\top A\mu.
\]
\end{proposition}
\begin{proof}
Write $Z=Y-\mu$, so $\E Z=0$. Expanding $Y=\mu+Z$ makes the two
cross terms have expectation zero. The remaining random term satisfies
\[
\E(Z^\top AZ)=\sum_{i,j}a_{ij}\E(Z_iZ_j)
=\sum_{i,j}a_{ij}\Cov(Y)_{ji}=\tr(A\Cov(Y)).
\]
Symmetry of $A$ is not needed for the identity.
\end{proof}

The \textbf{linear model} is
\[
Y=X\beta+\varepsilon,\qquad
\E\varepsilon=0,\qquad \Cov(\varepsilon)=\sigma^2I_n,
\quad\sigma^2>0.
\]
Here $X\in\R^{n\times p}$ is a fixed known design matrix and
$\beta\in\R^p$ is unknown. The response $Y$ and error $\varepsilon$ are
random; an observed response is denoted by $y$. Throughout Sections 1--5
assume $X$ has full column rank. For residual variance estimation require
$p<n$. An intercept is present precisely when the constant column
$\mathbf1=(1,\ldots,1)^\top$ belongs to $\operatorname{col}(X)$.
The model's mean vectors range over this column space. No normality
assumption is needed for the moment and optimality results below.

\section*{2. Fitted vectors, residuals, and variance}
Put
\[
K=X^\top X,\quad L=K^{-1}X^\top,\quad
H=XL=XK^{-1}X^\top,\quad M=I_n-H.
\]
The matrix $K$ is positive definite because
$a^\top Ka=\norm{Xa}^2>0$ when $a\ne0$.
The least-squares results from Unit 10 give
\[
\widehat\beta=LY,\qquad \widehat Y=HY,\qquad e=MY.
\]
The matrices $H$ and $M$ are the orthogonal projectors onto
$\operatorname{col}(X)$ and $\ker X^\top$, respectively. In particular
\[
H^\top=H=H^2,\quad M^\top=M=M^2,\quad
HX=X,\quad MX=0,\quad HM=MH=0.
\]
For instance, $H^2=XK^{-1}KK^{-1}X^\top=H$; symmetry follows from
symmetry of $K^{-1}$. Also $\operatorname{col}(H)=\operatorname{col}(X)$
because $H$ takes values there and fixes every $Xa$. Hence
\[
\operatorname{rank}H=\tr H=p,\qquad
\operatorname{rank}M=\tr M=n-p.
\]
The trace equality follows either from eigenvalues $0,1$ of a symmetric
idempotent matrix or from $\tr H=\tr(K^{-1}X^\top X)=p$.

\begin{theorem}[First and second moments of the fit]\label{thm:11:moments}
Under the linear-model assumptions,
\[
\begin{gathered}
\E\widehat\beta=\beta,\qquad \Cov(\widehat\beta)=\sigma^2K^{-1},\\
\E\widehat Y=X\beta,\qquad \Cov(\widehat Y)=\sigma^2H,\\
\E e=0,\qquad \Cov(e)=\sigma^2M,\qquad
\Cov(\widehat\beta,e)=0.
\end{gathered}
\]
If $p<n$, $s^2=\norm e^2/(n-p)$ is unbiased for $\sigma^2$.
\end{theorem}
\begin{proof}
Since $LX=I_p$, $\widehat\beta=\beta+L\varepsilon$, giving its mean and
covariance $\sigma^2LL^\top=\sigma^2K^{-1}KK^{-1}=\sigma^2K^{-1}$.
Likewise $\widehat Y=X\beta+H\varepsilon$ and $e=M\varepsilon$.
Idempotence yields $\Cov(\widehat Y)=\sigma^2H$ and
$\Cov(e)=\sigma^2M$. Their cross-covariance with the coefficients is
$\sigma^2LM=0$, because $X^\top M=0$.
Finally Proposition~\ref{prop:11:quadratic} gives
$\E\norm e^2=\E(\varepsilon^\top M\varepsilon)=\sigma^2\tr M
=\sigma^2(n-p)$. Divide by the positive number $n-p$.
\end{proof}

The model errors and residuals are different vectors:
$\varepsilon=Y-X\beta$, whereas $e=Y-X\widehat\beta=M\varepsilon$.
The residuals obey exact constraints $X^\top e=0$. Their variances are
$\sigma^2(1-h_{ii})$ and their off-diagonal covariances are
$-\sigma^2h_{ij}$. Zero cross-covariance of estimates and residuals
does not by itself establish independence.

\begin{example}[A fit, its geometry, and its uncertainty]\label{eg:11:fit}
Take
\[
t=(-3,-1,1,3)^\top,\quad X=(\mathbf1\ t),\quad y=(1,2,2,5)^\top.
\]
Here $\mathbf1^\top t=0$ and $t^\top t=20$, so
\[
K=\begin{pmatrix}4&0\\0&20\end{pmatrix},\qquad
X^\top y=\begin{pmatrix}10\\12\end{pmatrix},\qquad
\widehat\beta=\begin{pmatrix}5/2\\3/5\end{pmatrix}.
\]
The fitted and residual columns are
\[
\widehat y=\tfrac1{10}(7,19,31,43)^\top,\qquad
e=\tfrac1{10}(3,1,-11,7)^\top.
\]
Both $\mathbf1^\top e=0$ and $t^\top e=0$ hold by direct substitution.
The residual sum of squares is $\mathrm{SSE}=180/100=9/5$, and
$s^2=(9/5)/(4-2)=9/10$.
The covariance of the coefficient estimator is
$\sigma^2\operatorname{diag}(1/4,1/20)$; replacing $\sigma^2$ by $s^2$
estimates it by $\operatorname{diag}(9/40,9/200)$.

The projector has entries
$h_{ij}=1/4+t_it_j/20$, so its diagonal is
$(7/10,3/10,3/10,7/10)$. Thus the endpoint residuals have variance
$3\sigma^2/10$, and the middle residuals have variance $7\sigma^2/10$.
These are sampling variances of random residuals; they are not the squares
of the observed residual entries.
\end{example}

\begin{proposition}[Leverage]
For every $i$, $0\leqslant h_{ii}\leqslant1$, and $\sum_i h_{ii}=p$.
If an intercept is present, $h_{ii}\geqslant1/n$.
\end{proposition}
\begin{proof}
Let $d_i$ be the $i$th coordinate vector of $\R^n$. Since $H$ is a
projector, $h_{ii}=d_i^\top Hd_i=\norm{Hd_i}^2\geqslant0$.
The orthogonal decomposition $d_i=Hd_i+Md_i$ gives
$1=h_{ii}+\norm{Md_i}^2$, proving the upper bound.
Summing the diagonal gives $\tr H=p$.
If $\mathbf1\in\operatorname{col}(X)$, write
$J=\mathbf1\mathbf1^\top/n$. Then $HJ=JH=J$, and $H-J$ is symmetric
idempotent. Its diagonal is nonnegative, so $h_{ii}-1/n\geqslant0$.
\end{proof}

\section*{3. The Gauss--Markov theorem}
A \textbf{linear estimator} of $\beta$ has the form $BY$ for a fixed
$p\times n$ matrix $B$. It is \textbf{unbiased} for every $\beta$ exactly
when $BX=I_p$. Indeed $\E(BY)=BX\beta$, and equality to $\beta$ for all
$\beta$ is equivalent to equality of those matrices.
For symmetric matrices, $C\succeq D$ means $C-D$ is positive semidefinite.
This compares the variances of \emph{all} scalar linear combinations:
$a^\top Ca\geqslant a^\top Da$ for every $a$.

\begin{theorem}[Gauss--Markov]\label{thm:11:gm}
Among all linear unbiased estimators of $\beta$, $\widehat\beta=LY$ has
the smallest covariance matrix in the positive-semidefinite order.
More precisely, for any $B$ with $BX=I_p$,
\[
\Cov(BY)-\Cov(LY)=\sigma^2(B-L)(B-L)^\top\succeq0.
\]
Equality of the covariance matrices holds exactly when $B=L$.
\end{theorem}
\begin{proof}
Set $D=B-L$. Since $BX=LX=I_p$, $DX=0$.
Consequently $LD^\top=K^{-1}X^\top D^\top=0$ and $DL^\top=0$.
Expansion gives $BB^\top=LL^\top+DD^\top$.
Multiplying by $\sigma^2$ proves the covariance identity; for every $a$,
$a^\top DD^\top a=\norm{D^\top a}^2\geqslant0$.
If the covariance difference is zero, its trace is
$\sigma^2\sum_{i,j}d_{ij}^2=0$, forcing $D=0$.
Conversely $D=0$ plainly gives equality.
\end{proof}

For a single target $c^\top\beta$, equality of its variance means
$D^\top c=0$; it need not force every row of $D$ to vanish.
Unbiasedness is a condition for all parameter values. Choosing a matrix
using the observed responses would no longer give an estimator linear
in $Y$ in the sense used by this theorem.

\begin{example}[The excess variance is computable]
Continue Example~\ref{eg:11:fit}. Put $d=(1,-1,-1,1)^\top$.
Then $d^\top\mathbf1=d^\top t=0$. For any fixed real $a$, define
\[
\widetilde\beta_0=\widehat\beta_0+a\,d^\top Y,\qquad
\widetilde\beta_1=\widehat\beta_1.
\]
The extra term has mean $a d^\top X\beta=0$, so both estimators remain
unbiased. Since $\norm d^2=4$, the covariance difference is
$\sigma^2\operatorname{diag}(4a^2,0)$.
The slope variance is unchanged; the intercept variance increases strictly
unless $a=0$. With the observed data, $d^\top y=2$.
\end{example}

\section*{4. Removing nuisance columns and comparing models}
Partition a full-column-rank design as $X=(N\ Z)$, where $N$ contains
the nuisance columns. Write the coefficient column as $(\alpha^\top,\gamma^\top)^\top$.
Let $H_N=N(N^\top N)^{-1}N^\top$ and $M_N=I-H_N$.

\begin{theorem}[Partitioned regression]\label{thm:11:fwl}
The matrix $Z^\top M_NZ$ is positive definite, and
\[
\widehat\gamma=(Z^\top M_NZ)^{-1}Z^\top M_NY,\qquad
\widehat\alpha=(N^\top N)^{-1}N^\top(Y-Z\widehat\gamma).
\]
The full-model residual equals
$M_NY-M_NZ\widehat\gamma$.
\end{theorem}
\begin{proof}
If $M_NZv=0$, then $Zv\in\operatorname{col}(N)$, say $Zv=Nu$.
The relation $N(-u)+Zv=0$ and full column rank of $(N\ Z)$ give $v=0$.
Thus $v^\top Z^\top M_NZv=\norm{M_NZv}^2>0$ for $v\ne0$.

For a fixed $\gamma$, project $Y-Z\gamma$ onto $\operatorname{col}(N)$.
The unique minimizing $\alpha$ is
$(N^\top N)^{-1}N^\top(Y-Z\gamma)$, and the remaining residual is
$M_N(Y-Z\gamma)$. Minimizing its squared norm is now an ordinary
least-squares problem with response $M_NY$ and design $M_NZ$, of full
column rank as just proved. Its normal equations give the displayed
$\widehat\gamma$. Substitution yields $\widehat\alpha$ and the residual.
\end{proof}

\begin{example}[Centering a simple regression]
Let $N=\mathbf1$ and let $Z=z$ be a nonconstant observation column.
Then $M_Nz=z-\bar z\mathbf1$ and $M_Ny=y-\bar y\mathbf1$.
Partitioned regression gives
\[
\widehat\gamma=
\frac{\sum_i(z_i-\bar z)(y_i-\bar y)}{\sum_i(z_i-\bar z)^2},\qquad
\widehat\alpha=\bar y-\widehat\gamma\bar z.
\]
Centering removes the intercept before estimating the slope. In
Example~\ref{eg:11:fit}, $z=t$ has mean zero, the numerator is $12$ and
the denominator is $20$, recovering $\widehat\gamma=3/5$ and
$\widehat\alpha=5/2$. If all $z_i$ are equal, the denominator is zero:
the intercept and regressor columns are dependent, so separate intercept
and slope coefficients are not identifiable.
\end{example}

If $\mathcal S_0\subseteq\mathcal S_1$ are two model spaces, their orthogonal
projectors satisfy $H_1H_0=H_0H_1=H_0$. The difference $H_1-H_0$ is the
orthogonal projector onto $\mathcal S_1\cap\mathcal S_0^\perp$, since
its square equals itself and it fixes exactly that subspace. Therefore
\[
\norm{(I-H_0)y}^2
=\norm{(H_1-H_0)y}^2+\norm{(I-H_1)y}^2.
\]
Adding columns cannot increase the least-squares residual sum of squares.
It decreases it strictly exactly when $(H_1-H_0)y\ne0$.

When $\mathbf1\in\operatorname{col}(X)$, use $H_0=J=\mathbf1\mathbf1^\top/n$
to obtain
\[
\underbrace{\norm{y-\bar y\mathbf1}^2}_{\mathrm{SST}}
=\underbrace{\norm{\widehat y-\bar y\mathbf1}^2}_{\mathrm{SSR}}
+\underbrace{\norm{y-\widehat y}^2}_{\mathrm{SSE}}.
\]
For $\mathrm{SST}>0$, $R^2=1-\mathrm{SSE}/\mathrm{SST}$ lies in $[0,1]$.
For the running fit, $\mathrm{SST}=9$, $\mathrm{SSR}=36/5$ and
$\mathrm{SSE}=9/5$, giving $R^2=4/5$.
Without a constant column, this centered orthogonal decomposition need
not hold; the always-valid decomposition is
$\norm y^2=\norm{Hy}^2+\norm{My}^2$.

\section*{5. Generalized least squares}
Now suppose $\Cov(\varepsilon)=\sigma^2\Omega$, where $\Omega$ is known,
symmetric and positive definite. The means are unchanged. Let
$W=\Omega^{-1/2}$, the symmetric inverse square root supplied by the
spectral theorem. Then
\[
\widetilde Y=WY,\quad\widetilde X=WX,\quad
\widetilde\varepsilon=W\varepsilon,
\qquad\Cov(\widetilde\varepsilon)=\sigma^2I.
\]
This invertible change of observation coordinates is called whitening.

\begin{theorem}[Generalized least squares and Aitken's theorem]\label{thm:11:gls}
Let $K_\Omega=X^\top\Omega^{-1}X$. The unique minimizer of
$(Y-Xb)^\top\Omega^{-1}(Y-Xb)$ is
\[
\widehat\beta_\Omega=L_\Omega Y,\qquad
L_\Omega=K_\Omega^{-1}X^\top\Omega^{-1}.
\]
It is unbiased, with covariance $\sigma^2K_\Omega^{-1}$, and is best
among all linear unbiased estimators in the covariance order.
\end{theorem}
\begin{proof}
The objective is $\norm{WY-WXb}^2$ and $WX$ has full column rank.
Ordinary least squares on the whitened model gives the formula, its
unbiasedness and covariance. Alternatively, $L_\Omega X=I$ and
$L_\Omega\Omega L_\Omega^\top=K_\Omega^{-1}$ verify these directly.
For any $B$ with $BX=I$, set $D=B-L_\Omega$. Then $DX=0$ and
$L_\Omega\Omega D^\top=K_\Omega^{-1}X^\top D^\top=0$. Hence
\[
\Cov(BY)-\Cov(L_\Omega Y)=\sigma^2D\Omega D^\top\succeq0.
\]
Because $\Omega$ is positive definite, equality holds exactly when $D=0$.
\end{proof}

The original-coordinate fitted matrix
$H_\Omega=XL_\Omega$ is idempotent. It is generally not symmetric, but
$H_\Omega^\top\Omega^{-1}=\Omega^{-1}H_\Omega$.
The residual satisfies $X^\top\Omega^{-1}(Y-X\widehat\beta_\Omega)=0$:
orthogonality is measured by the weighted inner product
$\ip uv_\Omega=u^\top\Omega^{-1}v$.
In contrast, the ordinary estimator $LY$ is still unbiased, but its
actual covariance is the sandwich matrix
$\sigma^2K^{-1}X^\top\Omega XK^{-1}$.

\begin{example}[Two observations with unequal precision]
For $X=(1,1)^\top$ and $\Omega=\operatorname{diag}(1,4)$,
\[
\widehat\beta_\Omega=\tfrac45Y_1+\tfrac15Y_2,\qquad
\Var(\widehat\beta_\Omega)=\tfrac45\sigma^2.
\]
Ordinary least squares gives $(Y_1+Y_2)/2$ with variance $5\sigma^2/4$.
For $y=(1,3)^\top$, the estimates are $7/5$ and $2$, respectively.
The weighted residual is $(-2/5,8/5)^\top$, and its weighted sum is
$-2/5+(1/4)(8/5)=0$. The unweighted sum is $6/5$.
Here $H_\Omega=\begin{pmatrix}4/5&1/5\\4/5&1/5\end{pmatrix}$ is visibly
nonsymmetric, while it still fixes the constant columns and satisfies
$H_\Omega^2=H_\Omega$.
\end{example}

\section*{6. Constraints and nonidentifiable coefficients}
The following extensions use the same orthogonal-decomposition argument.
They also supply the geometric tools for the more substantial exercises.
Let $R$ be $q\times p$ of full row rank and impose $R\beta=r$.
The full-rank weighted objective splits as
\[
Q(b)=Q(\widehat\beta_\Omega)
+(b-\widehat\beta_\Omega)^\top K_\Omega(b-\widehat\beta_\Omega),
\]
because the weighted normal equations make the cross term zero.
Put $C=RK_\Omega^{-1}R^\top$; it is positive definite since
$R^\top a\ne0$ when $a\ne0$.

\begin{proposition}[Constrained least squares]\label{prop:11:constraint}
The unique constrained minimizer is
\[
\widehat\beta_R=\widehat\beta_\Omega
-K_\Omega^{-1}R^\top C^{-1}(R\widehat\beta_\Omega-r).
\]
Its increase in the objective is
$(R\widehat\beta_\Omega-r)^\top C^{-1}(R\widehat\beta_\Omega-r)$.
\end{proposition}
\begin{proof}
The formula gives $R\widehat\beta_R=r$, so it is feasible.
If another feasible vector is $\widehat\beta_R+v$, then $Rv=0$.
The cross term in its $K_\Omega$-distance from $\widehat\beta_\Omega$
is zero because
\[
v^\top K_\Omega(\widehat\beta_R-\widehat\beta_\Omega)
=-v^\top R^\top C^{-1}(R\widehat\beta_\Omega-r)=0.
\]
The additional objective is thus $v^\top K_\Omega v$, strictly positive
unless $v=0$. Substitution of the displayed displacement gives the
claimed increase, using $RK_\Omega^{-1}R^\top=C$.
\end{proof}

For the running fit, imposing slope zero gives $R=(0\ 1)$, $r=0$,
$C=1/20$. The constrained intercept remains $5/2$ and the increase
is $(3/5)^2/(1/20)=36/5$. Its SSE is $9/5+36/5=9$, the total centered
sum of squares, agreeing with the constant-only model.

If $X$ has rank $r<p$, the mean $X\beta$ is identifiable, but individual
coefficients may not be. Parameter columns $\beta$ and $\beta+v$ with
$v\in\ker X$ describe the same mean. A scalar target $c^\top\beta$ is
\textbf{estimable} if some linear estimator $a^\top Y$ is unbiased for
it for every $\beta$.

\begin{theorem}[Estimable targets]\label{thm:11:estimable}
The target $c^\top\beta$ is estimable exactly when
$c\in\operatorname{col}(X^\top)=(\ker X)^\perp$.
Under spherical covariance, $c^\top X^+Y$ is its best linear unbiased
estimator. The fitted vector is $XX^+Y$, and $\norm{(I-XX^+)Y}^2/(n-r)$
is unbiased for $\sigma^2$ when $r<n$.
\end{theorem}
\begin{proof}
Unbiasedness of $a^\top Y$ means $a^\top X\beta=c^\top\beta$ for every
$\beta$, or $X^\top a=c$. This proves the first equivalence; the
orthogonal-complement equality is the fundamental row/kernel relation.
For estimable $c$, $c^\top X^+X=c^\top$, since $X^+X$ projects onto
the row space. Thus $a_0=(X^+)^\top c$ satisfies $X^\top a_0=c$.
Any other solution is $a=a_0+d$ with $X^\top d=0$.
The SVD shows $a_0\in\operatorname{col}(X)$, so $a_0^\top d=0$.
Consequently $\Var(a^\top Y)=\sigma^2\norm{a_0}^2+\sigma^2\norm d^2$,
proving optimality. The last two claims follow from the projector
$H=XX^+$ of rank $r$ and the earlier moment proof, with $p$ replaced
by $r$.
\end{proof}

For $X=(\mathbf1\ t\ 2t)$ with the centered nonzero column $t$ from
Example~\ref{eg:11:fit}, the mean is
$\beta_0\mathbf1+(\beta_1+2\beta_2)t$.
The estimable targets are exactly those with coefficient vector
$(c_0,c_1,2c_1)^\top$. The separate slopes are not estimable.
All least-squares coefficient columns satisfy
$b_0=5/2$, $b_1+2b_2=3/5$. Their minimum-norm choice is
$(5/2,3/25,6/25)^\top$, obtained by projecting $(b_1,b_2)$ onto
the direction $(1,2)$. Fitted values and SSE are unchanged by this
redundant parameterization.
